% ============================================================
%  TinyMPC + Differential Flatness: Complete Cheat Sheet
%  Jianghan Zhang, 2026
%  5-variant hierarchy: Online Relin → Flat Template → Barrier → Analytic → Gain Schedule
% ============================================================
\documentclass[11pt]{article}

% ── packages ──
\usepackage[margin=1in]{geometry}
\usepackage{amsmath,amssymb,amsthm}
\usepackage{mathtools}
\usepackage{algorithm2e}
\usepackage{booktabs}
\usepackage{enumitem}
\usepackage{hyperref}
\usepackage{cleveref}
\usepackage{listings}
\usepackage{xcolor}

% ── listing style ──
\lstset{
  language=Python,
  basicstyle=\ttfamily\small,
  keywordstyle=\color{blue!70!black}\bfseries,
  commentstyle=\color{green!50!black}\itshape,
  stringstyle=\color{red!60!black},
  numbers=left,
  numberstyle=\tiny\color{gray},
  numbersep=5pt,
  frame=single,
  breaklines=true,
  columns=fullflexible,
  showstringspaces=false,
  tabsize=4,
}

% ── theorem environments ──
\newtheorem{theorem}{Theorem}[section]
\newtheorem{proposition}[theorem]{Proposition}
\newtheorem{lemma}[theorem]{Lemma}
\newtheorem{corollary}[theorem]{Corollary}
\newtheorem{definition}[theorem]{Definition}
\newtheorem{remark}[theorem]{Remark}
\newtheorem{assumption}[theorem]{Assumption}

% ── notation ──
\newcommand{\R}{\mathbb{R}}
\newcommand{\SO}{\mathrm{SO}}
\newcommand{\SE}{\mathrm{SE}}
\newcommand{\rank}{\mathrm{rank}}
\newcommand{\diag}{\mathrm{diag}}
\newcommand{\clip}{\mathrm{clip}}
\newcommand{\prox}{\mathrm{prox}}
\newcommand{\spec}{\mathrm{spec}}
\newcommand{\tr}{\mathrm{tr}}
\newcommand{\ip}[2]{\langle #1,\, #2 \rangle}
\newcommand{\norm}[1]{\|#1\|}
\newcommand{\trans}{\mathsf{T}}
\DeclareMathOperator*{\argmin}{arg\,min}
\DeclareMathOperator{\vect}{vec}

\title{\textbf{TinyMPC + Differential Flatness}:\\
Complete Cheat Sheet (5-Variant Hierarchy)}
\author{Jianghan Zhang}
\date{February 2026}

\begin{document}
\maketitle

\begin{abstract}
TinyMPC solves linear MPC via ADMM with offline-cached
Riccati matrices, achieving microcontroller-grade speed
at the cost of a fixed linearisation point.
We observe that for differentially flat systems
the linearisation is a smooth function of
the flat output and its derivatives---a
low-dimensional parameter $\alpha\in\R^d$
(e.g.\ $d=4$ for a quadrotor).
This gives a family of controllers forming a
\textbf{complexity cascade}:
\begin{center}
\small
Online Relin $\xrightarrow{\text{cache}}$
Flat Template $\xrightarrow{K\to 0}$
Barrier Templ.\ $\xrightarrow{\mu\to 0}$
Analytic Templ.\ $\xrightarrow{\text{1st order}}$
Gain Schedule
\end{center}
Moving down the cascade trades constraint fidelity
for real-time guarantees, from $O(n^3)$ online NMPC
to ${\sim}260$-operation deterministic feedback.
\end{abstract}

\tableofcontents
\newpage

% ============================================================
\section{TinyMPC Formulation}\label{sec:tinympc}
% ============================================================

We follow the formulation of \cite{tinympc}.

\subsection{Problem}

Given discrete-time linear dynamics
\begin{equation}\label{eq:dynamics}
  x_{k+1} = A\,x_k + B\,u_k,
  \qquad x_k \in \R^n,\; u_k \in \R^m,
\end{equation}
TinyMPC solves the finite-horizon constrained LQR:
\begin{equation}\label{eq:mpc}
\begin{aligned}
  \min_{x,u} \quad
    & \sum_{k=0}^{N-1}
      \Bigl(\tfrac{1}{2}x_k^\trans Q\, x_k
      + \tfrac{1}{2}u_k^\trans R\, u_k
      + q^\trans x_k + r^\trans u_k\Bigr)
      + \tfrac{1}{2}x_N^\trans Q_f\, x_N
      + q_f^\trans x_N \\
  \text{s.t.} \quad
    & x_{k+1} = A\,x_k + B\,u_k,
      \qquad k = 0,\ldots,N{-}1, \\
    & x_{\min} \le x_k \le x_{\max}, \\
    & u_{\min} \le u_k \le u_{\max},
\end{aligned}
\end{equation}
with $Q, Q_f \succeq 0$, $R \succ 0$.

\subsection{ADMM solver}\label{ssec:admm}

Introduce slack variables $v = (v_x, v_u)$
to decouple dynamics from box constraints.
The augmented Lagrangian with penalty $\rho > 0$ is
\begin{equation}\label{eq:auglag}
  \mathcal{L}_\rho
  = J(x,u) + \sum_k \Bigl(
    \lambda_{x,k}^\trans(x_k - v_{x,k})
    + \lambda_{u,k}^\trans(u_k - v_{u,k})
    + \tfrac{\rho}{2}\norm{x_k - v_{x,k}}^2
    + \tfrac{\rho}{2}\norm{u_k - v_{u,k}}^2
  \Bigr).
\end{equation}

Each ADMM iteration:
\begin{enumerate}[label=(\roman*), leftmargin=8mm]
  \item \textbf{Primal update}
    (Riccati backward pass).
    Solve the unconstrained LQR with shifted cost:
    \begin{equation}\label{eq:riccati}
    \begin{aligned}
      & P_N = Q_f + \rho\, I, \\
      & K_k = (R + \rho\, I + B^\trans P_{k+1}\, B)^{-1}
              B^\trans P_{k+1}\, A, \\
      & P_k = Q + \rho\, I + A^\trans P_{k+1}\, A
              - A^\trans P_{k+1}\, B\, K_k.
    \end{aligned}
    \end{equation}
    The gain $K_k$ and the affine term
    $d_k = (R + \rho\, I + B^\trans P_{k+1}\, B)^{-1}
    (B^\trans p_{k+1} + r + \rho\,(v_{u,k} - \lambda_{u,k}/\rho))$
    produce the primal trajectory via forward rollout:
    $u_k = -K_k\, x_k - d_k$,
    $x_{k+1} = A\,x_k + B\,u_k$.

  \item \textbf{Slack update}
    (element-wise projection):
    \begin{equation}\label{eq:slack}
      v_{x,k} = \clip(x_k + \lambda_{x,k}/\rho,\;
                       x_{\min},\; x_{\max}),
      \qquad
      v_{u,k} = \clip(u_k + \lambda_{u,k}/\rho,\;
                       u_{\min},\; u_{\max}).
    \end{equation}

  \item \textbf{Dual update}:
    \begin{equation}\label{eq:dual}
      \lambda_{x,k} \mathrel{+}= \rho\,(x_k - v_{x,k}),
      \qquad
      \lambda_{u,k} \mathrel{+}= \rho\,(u_k - v_{u,k}).
    \end{equation}
\end{enumerate}

\begin{remark}[Offline cache]\label{rem:cache}
When $A, B, Q, R, \rho$ are fixed,
the Riccati matrices $\{P_k, K_k\}$ converge
to steady-state values $P_\infty, K_\infty$
as $N \to \infty$.
TinyMPC precomputes and caches $P_\infty, K_\infty$
and the matrix
$C_\infty := (R + \rho\, I + B^\trans P_\infty\, B)^{-1}$
offline.
Online cost: one forward pass (matrix--vector products)
$+$ two element-wise projections $+$ one dual update.
No matrix inversions at runtime.
\end{remark}

\subsection{Limitation: fixed linearisation}

For a nonlinear system $x_{k+1} = f(x_k, u_k)$,
TinyMPC uses a single linearisation
\begin{equation}\label{eq:fixed-lin}
  A = \frac{\partial f}{\partial x}\bigg|_{(\bar x, \bar u)},
  \qquad
  B = \frac{\partial f}{\partial u}\bigg|_{(\bar x, \bar u)},
\end{equation}
at a fixed nominal $(\bar x, \bar u)$
(typically hover for a quadrotor).
The cached matrices are valid only near $(\bar x, \bar u)$.
Deviation $\to$ model mismatch $\to$ degraded or
lost feasibility.

% ============================================================
\section{Differential Flatness}\label{sec:flatness}
% ============================================================

\begin{definition}[Differentially flat system]
\label{def:flat}
A system $\dot x = f(x, u)$ with state $x\in\R^n$,
input $u\in\R^m$ is \textbf{differentially flat}
if there exist \textbf{flat outputs}
$\sigma = h(x, u, \dot u, \ldots, u^{(r)}) \in \R^m$
such that:
\begin{enumerate}[label=(\alph*), leftmargin=8mm]
  \item $x = \varphi(\sigma, \dot\sigma, \ldots,
    \sigma^{(s)})$
    \quad (state recovery),
  \item $u = \psi(\sigma, \dot\sigma, \ldots,
    \sigma^{(s+1)})$
    \quad (input recovery),
\end{enumerate}
for smooth maps $\varphi, \psi$ and some finite $s$.
\end{definition}

\begin{remark}[Quadrotor]
For a quadrotor with state
$(p, v, R, \omega) \in \R^3 \times \R^3 \times \SO(3)
\times \R^3$
and inputs $(T, \tau) \in \R \times \R^3$
(thrust + body torques),
the flat outputs are
\begin{equation}\label{eq:flat-quad}
  \sigma = (x, y, z, \psi) \in \R^3 \times S^1
  \;\approx\; \R^4.
\end{equation}
Recovery:
position $p = \sigma_{1:3}$,
velocity $v = \dot\sigma_{1:3}$,
thrust direction
$e_3^{\mathcal{B}} = (m\ddot\sigma_{1:3} + mg\,e_3)/
\norm{m\ddot\sigma_{1:3} + mg\,e_3}$
(determines roll and pitch from $\ddot\sigma$),
yaw from $\sigma_4 = \psi$.
The full rotation $R$ and angular velocity $\omega$
follow from $(\ddot\sigma, \dddot\sigma, \psi, \dot\psi)$.
Thrust $T = \norm{m\ddot\sigma_{1:3} + mg\,e_3}$.
\end{remark}

\subsection{Flat equilibrium manifold}

Given a flat output trajectory
$\sigma(t)$ and its derivatives,
the flatness maps produce a
\textbf{dynamically consistent} pair:
\begin{equation}\label{eq:flat-eq}
  x^*(\alpha) = \varphi(\alpha),
  \qquad
  u^*(\alpha) = \psi(\alpha),
  \qquad
  \alpha := (\sigma, \dot\sigma, \ldots, \sigma^{(s+1)}).
\end{equation}
This pair satisfies the dynamics exactly:
$f(x^*(\alpha), u^*(\alpha)) = \dot x^*(\alpha)$.

\begin{definition}[Flat parameter]\label{def:flat-param}
For a given flat system,
the \textbf{flat parameter} is the minimal
sub-tuple of $\alpha$ on which the
linearisation $(A, B)$ depends.
Denote it $\alpha \in \R^d$.
\end{definition}

\begin{remark}[Quadrotor flat parameter]\label{rem:quad-param}
For a quadrotor, the Jacobians
$\partial f/\partial x$ and $\partial f/\partial u$
evaluated at $(x^*, u^*)$ depend only on:
\begin{equation}\label{eq:quad-alpha}
  \alpha = (\ddot\sigma_{1:3},\, \psi)
  = (a_x, a_y, a_z, \psi) \in \R^4,
\end{equation}
where $(a_x, a_y, a_z)$ is the commanded acceleration
(determining the equilibrium attitude)
and $\psi$ is the yaw angle.
At hover: $\alpha_0 = (0, 0, 0, 0)$.
\end{remark}

% ============================================================
\section{Variant 1: Online Relinearization}\label{sec:online-relin}
% ============================================================

The most accurate but most expensive variant.
At each control step:
\begin{enumerate}[leftmargin=6mm]
  \item Compute $\alpha$ from the trajectory reference.
  \item Evaluate $(A(\alpha), B(\alpha))$ via autograd or
    finite-difference Jacobians.
  \item Solve a fresh DARE: $O(n^3)$.
  \item Run $K$ ADMM iterations with the new cache: $K \cdot O(n^2)$.
\end{enumerate}

\textbf{Online cost:} $O(n^3 + K\,n^2)$ per step.

This is the ``gold standard'' that all other variants
approximate.
The per-step DARE solve is the bottleneck
(${\sim}22$\,s total for 200 figure-eight steps in Python).

% ============================================================
\section{Variant 2: Flat Template}\label{sec:template}
% ============================================================

\subsection{Parameterised linearisation}

Linearise $f$ at the flat equilibrium
$(x^*(\alpha), u^*(\alpha))$:
\begin{equation}\label{eq:param-lin}
  A(\alpha) = \frac{\partial f}{\partial x}
    \bigg|_{(x^*(\alpha),\, u^*(\alpha))},
  \qquad
  B(\alpha) = \frac{\partial f}{\partial u}
    \bigg|_{(x^*(\alpha),\, u^*(\alpha))}.
\end{equation}

\begin{proposition}[Smoothness]\label{prop:smooth}
If $f, \varphi, \psi$ are $C^2$, then
$\alpha \mapsto (A(\alpha), B(\alpha))$
is $C^1$ on $\R^d$.
\end{proposition}
\begin{proof}
Composition of $C^2$ maps
($f$, $\varphi$, $\psi$, and the Jacobian operator)
is $C^1$.
\end{proof}

\subsection{Riccati cache family}

For each $\alpha$, the steady-state Riccati equation
with $A(\alpha), B(\alpha)$ and fixed $Q, R, \rho$ gives:
\begin{equation}\label{eq:riccati-alpha}
\begin{aligned}
  P_\infty(\alpha)
  &= Q + \rho\, I + A(\alpha)^\trans P_\infty(\alpha)\, A(\alpha)
  \\
  &\quad - A(\alpha)^\trans P_\infty(\alpha)\, B(\alpha)\,
    K_\infty(\alpha), \\[4pt]
  K_\infty(\alpha)
  &= \bigl(R + \rho\, I + B(\alpha)^\trans
    P_\infty(\alpha)\, B(\alpha)\bigr)^{-1}
    B(\alpha)^\trans P_\infty(\alpha)\, A(\alpha), \\[4pt]
  C_\infty(\alpha)
  &= \bigl(R + \rho\, I + B(\alpha)^\trans
    P_\infty(\alpha)\, B(\alpha)\bigr)^{-1}.
\end{aligned}
\end{equation}

\begin{proposition}[Smooth cache]\label{prop:smooth-cache}
If $(A(\alpha), B(\alpha))$ is stabilisable for all
$\alpha$ in a compact set $\mathcal{A} \subset \R^d$,
then $\alpha \mapsto (P_\infty(\alpha), K_\infty(\alpha),
C_\infty(\alpha))$ is $C^1$ on $\mathcal{A}$.
\end{proposition}

\subsection{The template}

\begin{definition}[Flat template]\label{def:template}
A \textbf{flat template} for a differentially flat
system with cost $(Q, R, \rho)$ is a map
\begin{equation}\label{eq:template}
  \mathcal{T} : \R^d \to
  \R^{n\times n} \times \R^{m\times n} \times \R^{m\times m},
  \qquad
  \alpha \mapsto
  \bigl(P_\infty(\alpha),\, K_\infty(\alpha),\,
  C_\infty(\alpha)\bigr),
\end{equation}
defined by the DARE \eqref{eq:riccati-alpha}.
\end{definition}

\noindent
\textbf{Offline}: evaluate $\mathcal{T}$
on a grid $\{\alpha_i\}_{i=1}^G \subset \mathcal{A}$
and store the cache triples.

\noindent
\textbf{Online}: given current flat output derivatives,
compute $\alpha_{\mathrm{now}}$,
look up (or interpolate) the nearest cache,
and run TinyMPC's ADMM
with $(A(\alpha_{\mathrm{now}}),\,
B(\alpha_{\mathrm{now}}),\,
K_\infty(\alpha_{\mathrm{now}}),\,
C_\infty(\alpha_{\mathrm{now}}))$.

\textbf{Online cost:} $K \cdot O(n^2)$ (same as TinyMPC per iteration).

\begin{remark}[Storage]\label{rem:storage}
For a quadrotor ($n=12, m=4, d=4$):
each cache triple requires
$12^2 + 4 \times 12 + 4^2 = 208$ floats
$= 832$ bytes (float32).
A $10^4$-point grid in $\R^4$
costs $\approx 8$ MB---within microcontroller flash.
\end{remark}

% ============================================================
\section{Flat Template Algorithm}\label{sec:algorithm}
% ============================================================

\begin{algorithm}[H]
\SetAlgoLined
\caption{Flat Template MPC}\label{alg:flat-mpc}
\KwIn{Current state $x_0$;\;
  flat output derivatives
  $(\sigma, \dot\sigma, \ldots)$ from planner or estimator;\;
  template $\mathcal{T}$;\;
  ADMM iterations $K$;\;
  penalty $\rho > 0$;\;
  warm-start $(\bar v, \bar\lambda)$.}
\KwOut{Control $u_0^* \in \R^m$.}

\medskip

\textbf{Step 0: Flat parameter.}\;
$\alpha \gets (\ddot\sigma_{1:3},\, \psi)$
\tcp*{quadrotor case}

\textbf{Step 1: Template lookup.}\;
$(A, B) \gets (A(\alpha),\, B(\alpha))$
\tcp*{evaluate or interpolate}
$(P, K, C) \gets \mathcal{T}(\alpha)$
\tcp*{cached Riccati}

\textbf{Step 2: Reference.}\;
$x^* \gets \varphi(\alpha)$,\;
$u^* \gets \psi(\alpha)$
\tcp*{flat equilibrium}

Shift coordinates:
$\tilde x_0 \gets x_0 - x^*$,\;
update linear cost terms
$q \gets Q\, x^*$, $r \gets R\, u^*$\;

\medskip

\textbf{Step 3: ADMM (identical to TinyMPC).}\;
\For{$j = 1, \ldots, K$}{
  \textbf{Primal:}
  Forward rollout with cached $K, C$:\;
  $\tilde u_k = -K\,\tilde x_k - C\,(B^\trans p_{k+1}
  + r + \rho(v_{u,k} - \lambda_{u,k}/\rho))$,\;
  $\tilde x_{k+1} = A\,\tilde x_k + B\,\tilde u_k$\;

  \textbf{Slack:}
  $v_{x,k} = \clip(\tilde x_k + x^* + \lambda_{x,k}/\rho,\;
    x_{\min},\; x_{\max}) - x^*$\;
  $v_{u,k} = \clip(\tilde u_k + u^* + \lambda_{u,k}/\rho,\;
    u_{\min},\; u_{\max}) - u^*$\;

  \textbf{Dual:}
  $\lambda_{x,k} \mathrel{+}= \rho\,(\tilde x_k - v_{x,k})$,\;
  $\lambda_{u,k} \mathrel{+}= \rho\,(\tilde u_k - v_{u,k})$\;
}

\Return $u_0^* = \tilde u_0 + u^*$\;
\end{algorithm}

% ============================================================
\section{Uniform Tunnel Guarantee}\label{sec:tunnel}
% ============================================================

We connect the flat template to the tunnel theory
\cite{zhang2026ttt}.

\subsection{Tunnel in original coordinates}

\begin{definition}[Tunnel radius]\label{def:tunnel}
For a fixed linearisation $(A, B)$ at nominal
$(\bar x, \bar u)$, the \textbf{tunnel radius} is
\begin{equation}\label{eq:tunnel-radius}
  r(A, B) = \frac{2\,\alpha(A, B)}{L},
\end{equation}
where $\alpha(A, B) = 1 - \bar\rho(A - B\,K_\infty)$
is the spectral gap of the closed-loop operator
(using $\bar\rho$ for spectral radius to avoid
overloading with ADMM penalty $\rho$)
and $L$ is the Lipschitz constant of $\nabla^2 f$.
\end{definition}

\begin{assumption}[Uniform stabilisability]
\label{ass:uniform}
There exists $\bar\alpha > 0$ such that
for all $\alpha \in \mathcal{A}$:
\begin{enumerate}[label=(\roman*), leftmargin=8mm]
  \item $(A(\alpha), B(\alpha))$ is stabilisable,
  \item $\alpha(A(\alpha), B(\alpha)) \ge \bar\alpha > 0$.
\end{enumerate}
\end{assumption}

\begin{theorem}[Uniform tunnel]\label{thm:uniform}
Under Assumption~\ref{ass:uniform},
the flat template MPC
(Algorithm~\ref{alg:flat-mpc})
has a \textbf{uniform} tunnel radius
\begin{equation}\label{eq:uniform-tunnel}
  r_{\min} = \inf_{\alpha \in \mathcal{A}}\,
  r(\alpha) \ge \frac{2\,\bar\alpha}{L} > 0
\end{equation}
at every flat equilibrium in $\mathcal{A}$.
The ADMM iterates contract at rate
$\le 1 - \bar\alpha$ uniformly:
\begin{equation}\label{eq:uniform-contract}
  \norm{w^{(j+1)} - w^*(\alpha)}
  \le (1 - \bar\alpha)\,
  \norm{w^{(j)} - w^*(\alpha)}
\end{equation}
for all $\alpha \in \mathcal{A}$
and all $w^{(0)}$ with
$\norm{w^{(0)} - w^*(\alpha)} < r_{\min}$,
where $w = (x, u, v, \lambda)$ is the full
primal-dual variable.
\end{theorem}

\begin{corollary}[Full flat-reachable envelope]
\label{cor:coverage}
The union of all $\alpha$-tunnels covers
a neighbourhood of the entire
flat equilibrium manifold.
For a quadrotor with $\mathcal{A}$ spanning
accelerations up to $a_{\max}$ and all yaw angles,
this includes all attitudes reachable
by $\norm{a} \le a_{\max}$---far beyond hover.
\end{corollary}

% ============================================================
\section{Variant 3: Analytic Template (Zero-Iteration MPC)}
\label{sec:analytic}
% ============================================================

The flat template (\S\ref{sec:template})
eliminates the fixed-nominal limitation
but still runs $K$ ADMM iterations online.
We now eliminate the iterations entirely.

\subsection{Key observation}

When box constraints are inactive,
the ADMM solution after $K\to\infty$ iterations
reduces to the unconstrained LQR:
\begin{equation}\label{eq:zero-iter}
  u_0^* = u^*(\alpha) - K_\infty(\alpha)\,
  (x_0 - x^*(\alpha)).
\end{equation}
This is a \textbf{single matrix--vector product}---no
iteration, no slack update, no dual update.

\subsection{Quadrotor: explicit $A(\alpha), B(\alpha)$}

For a quadrotor with state
$x = (p, \phi, \theta, \psi,
v, \omega)^\trans
\in \R^{12}$
and input $u = (u_1, u_2, u_3, u_4)^\trans \in \R^4$
(individual motor thrusts),
the flat parameter is
$\alpha = (a_x, a_y, a_z, \psi)$.

The equilibrium attitude follows from the
thrust direction:
\begin{equation}\label{eq:eq-attitude}
  T^* = m\norm{a + g\,e_3},
  \qquad
  \phi^* = \arcsin\!\Bigl(
    \frac{m(a_x \sin\psi - a_y \cos\psi)}{T^*}
  \Bigr),
  \qquad
  \theta^* = \arctan\!\Bigl(
    \frac{a_x \cos\psi + a_y \sin\psi}{a_z + g}
  \Bigr).
\end{equation}

The linearised dynamics at this equilibrium have
the block structure:
\begin{equation}\label{eq:A-explicit}
  A_c(\alpha) =
  \begin{pmatrix}
    0 & 0 & I_3 & 0 \\
    0 & 0 & 0 & W(\phi^*, \theta^*) \\
    0 & G(\alpha) & 0 & 0 \\
    0 & 0 & 0 & 0
  \end{pmatrix}\!,
  \quad
  B_c(\alpha) =
  \begin{pmatrix}
    0 \\ 0 \\ b(\alpha) \cdot \mathbf{1}^\trans \\ J^{-1} M_{\text{mix}}
  \end{pmatrix}\!,
\end{equation}
where $G(\alpha) \in \R^{3\times 3}$ encodes
the gravity coupling,
$W(\phi^*,\theta^*)$ is the Euler kinematic map,
$b(\alpha) = k_t \cdot e_3^{\mathcal{B}}(\alpha)/m$,
and $J^{-1}M_{\text{mix}}$ is the torque allocation.
Discretization: $A_d = I + A_c \Delta t$,
$B_d = B_c \Delta t$.

\subsection{Perturbation expansion of $K_\infty$}

Write $\alpha = \alpha_0 + \delta\alpha$
where $\alpha_0 = (0, 0, 0, 0)$ is hover.
The DARE perturbation theory gives:
\begin{equation}\label{eq:dP}
  P_\infty(\alpha)
  = P_0 + \delta P + O(\norm{\delta\alpha}^2),
\end{equation}
where $\delta P$ solves the discrete Lyapunov equation
\begin{equation}\label{eq:lyap}
  \delta P
  = A_{\mathrm{cl},0}^\trans\, \delta P\,
    A_{\mathrm{cl},0}
  + S(\delta A, \delta B, P_0, K_0),
\end{equation}
with $A_{\mathrm{cl},0} = A_0 - B_0\, K_0$.

\begin{proposition}[Analytic gain]
\label{prop:analytic-gain}
The first-order analytic gain is
\begin{equation}\label{eq:K-analytic}
  K(\alpha) = K_0
  + \sum_{i=1}^{d} \delta\alpha_i\, K_i
  + O(\norm{\delta\alpha}^2),
\end{equation}
where each $K_i \in \R^{m \times n}$
is a constant matrix given by
\begin{equation}\label{eq:Ki}
  K_i = C_0\bigl(
    \delta B_i^\trans P_0\, A_0
    + B_0^\trans P_i\, A_0
    + B_0^\trans P_0\, \delta A_i
    - \delta B_i^\trans P_0\, B_0\, K_0
    - B_0^\trans P_i\, B_0\, K_0
  \bigr),
\end{equation}
with $C_0 = (R + B_0^\trans P_0\, B_0)^{-1}$.
All quantities are constant.
\end{proposition}

\subsection{The zero-iteration control law}

Combining \eqref{eq:zero-iter}
and \eqref{eq:K-analytic}:

\begin{equation}\label{eq:brutal}
\boxed{
  u = u^*(\alpha)
  - \Bigl(K_0 + \sum_{i=1}^{d}
    \delta\alpha_i\, K_i\Bigr)
  \bigl(x_0 - x^*(\alpha)\bigr).
}
\end{equation}

\noindent
\textbf{Online cost:} $O(d\,mn)$.
For a quadrotor: $O(4 \cdot 4 \cdot 12) = O(192)$
multiply-accumulate operations.
With equilibrium computation: $\approx 260$ total ops.

\begin{remark}[Constraints]\label{rem:constraints}
When constraints are active, clamp the output:
$u = \clip(u^*(\alpha) - K(\alpha)(x_0 - x^*(\alpha)),
\; u_{\min},\; u_{\max})$.
This is a zeroth-order ADMM approximation.
\end{remark}

% ============================================================
\section{Variant 4: Barrier-Augmented Template}
\label{sec:barrier}
% ============================================================

The Analytic Template handles constraints only by
clipping---the gain $K(\alpha)$ is unaware of limits.
We now show that log-barrier functions can be
absorbed into the DARE cache, yielding a zero-iteration
controller that is inherently constraint-aware.

\subsection{Interior-point formulation}

For box constraints $u_{\min} \le u \le u_{\max}$,
the log-barrier method replaces these with a smooth penalty:
\begin{equation}\label{eq:barrier}
  \mathcal{B}_\mu(u) = -\mu \sum_{j=1}^{m}
    \bigl[\log(u_j - u_{\min,j})
          + \log(u_{\max,j} - u_j)\bigr],
\end{equation}
where $\mu > 0$ is the barrier parameter.

\subsection{Input barrier Hessian}

At the flat equilibrium $u^*(\alpha)$, the barrier
Hessian is a known diagonal matrix:
\begin{equation}\label{eq:barrier_hess}
  H_\mu(\alpha) = \mu \cdot \diag\!\left(
    \frac{1}{(u^*_j - u_{\min,j})^2}
    + \frac{1}{(u_{\max,j} - u^*_j)^2}
  \right).
\end{equation}
Large when $u^*(\alpha)$ is near a bound;
small when far from all bounds.

\subsection{State barrier Hessian}

For state constraints $x_{\min} \le x \le x_{\max}$
(e.g., position bounds $|p_i| \le L$):
\begin{equation}\label{eq:barrier_hess_x}
  H_{x,\mu}(\alpha) = \mu \cdot \diag\!\left(
    \frac{1}{(x^*_i - x_{\min,i})^2}
    + \frac{1}{(x_{\max,i} - x^*_i)^2}
  \right),
\end{equation}
applied only to constrained state dimensions
(unconstrained dimensions contribute zero).
In error coordinates ($x^*_{\mathrm{pos}} = 0$),
this depends on the reference-to-boundary distance.

\subsection{Barrier-augmented DARE}

Define barrier-augmented costs:
\begin{align}
  R_\mu(\alpha) &= R + H_\mu(\alpha), \label{eq:R_barrier} \\
  Q_\mu(\alpha) &= Q_{\mathrm{lqr}} + H_{x,\mu}(\alpha), \label{eq:Q_barrier}
\end{align}
where $Q_{\mathrm{lqr}}$ is the terminal cost from
standard discrete LQR.

Solving the augmented DARE:
\begin{equation}\label{eq:dare_barrier}
  P_\mu = Q_{\mu,\rho} + A_d^\top P_\mu \bigl(
    A_d - B_d \, G_\mu^{-1} \,
    B_d^\top P_\mu A_d \bigr),
\end{equation}
where $Q_{\mu,\rho} = Q_\mu + \rho I$ and
$G_\mu = R_\mu + \rho I + B_d^\top P_\mu B_d$,
yields the barrier-augmented gain
$K_\mu(\alpha) = G_\mu^{-1} B_d^\top P_\mu A_d$.

The zero-iteration barrier control law:
\begin{equation}\label{eq:barrier_control}
\boxed{
  u = u^*(\alpha)
  - \Bigl(K_{\mu,0} + \sum_{i=1}^{d}
    \delta\alpha_i\, \frac{\partial K_\mu}{\partial \alpha_i}\Bigr)
  \bigl(x_0 - x^*(\alpha)\bigr).
}
\end{equation}

\textbf{Online cost:} $\sim$260 ops (identical to Analytic Template).

\textbf{Storage:} $5 \times 4 \times 12 = 240$ floats (identical).

\subsection{$\alpha$-dependent conservatism}

At hover, the four motors produce equal thrust at midrange
and $H_\mu$ is moderate.
At aggressive maneuvers (large $\|a\|$), some motors
approach saturation and $H_\mu$ grows in those channels,
automatically increasing conservatism.
The scheduling variable $\alpha$ simultaneously parameterizes
the linearization, the equilibrium, \emph{and} the
constraint tightness.

\subsection{Optional: IPM refinement}

For tighter constraint satisfaction,
$K_{\mathrm{IPM}}$ Newton steps can be run online:
\begin{enumerate}[leftmargin=6mm]
  \item Start from the zero-iteration output $u^{(0)}$.
  \item Recompute $H_\mu$ at current $u^{(k-1)}$
    (not at equilibrium).
  \item Solve one Riccati backward pass with updated
    $R_\mu \to \delta u$.
  \item $u^{(k)} = u^{(k-1)} + \eta\, \delta u$
    (Newton step).
\end{enumerate}
Each step costs $O(n^2)$ (same as one ADMM iteration).
Unlike ADMM (linear convergence), Newton--IPM converges
\emph{quadratically}: 2--3 iterations typically suffice.

\subsection{Limitation: reference outside feasible region}

The barrier Hessian penalizes error-state excursions
from the reference. If the reference itself exceeds
the constraint boundary
(e.g., a circular trajectory outside a square bound),
no error-state cost modification can prevent violation.
Only the ADMM projection (which maps absolute bounds
to per-step error-state bounds) can handle this case.

\textbf{Key result:}
In the circle-in-square experiment,
Barrier = Analytic = 88.4\,mm max violation,
because $R - L = 88$\,mm is the geometric overshoot
of the reference.

% ============================================================
\section{Barrier-Augmented Algorithm}\label{sec:barrier-alg}
% ============================================================

\begin{algorithm}[H]
\SetAlgoLined
\caption{Barrier-Augmented Template MPC}\label{alg:barrier}
\KwIn{Constraint bounds $u_{\min}, u_{\max}$,
  $x_{\min}, x_{\max}$ (optional),
  barrier $\mu$, penalty $\rho$.}
\KwOut{$K_{\mu,0}$,
  $\{\partial K_\mu / \partial \alpha_i\}_{i=1}^4$,
  $\alpha_0$.}

\medskip
\textbf{Offline:}\\
$\alpha_0 \gets (0,0,0,0)$ \tcp*{hover}
$(x^*, u^*) \gets \Phi(\alpha_0)$
\tcp*{flat equilibrium}
$H_\mu \gets$ Eq.~\eqref{eq:barrier_hess}
from $u^*$ and input bounds\\
$H_{x,\mu} \gets$ Eq.~\eqref{eq:barrier_hess_x}
from $x^*$ and state bounds\\
$R_\mu \gets R + H_\mu$;\quad
$Q_\mu \gets Q_{\mathrm{lqr}} + H_{x,\mu}$\\
Solve DARE~\eqref{eq:dare_barrier}
with $R_\mu, Q_\mu \to K_{\mu,0}$\\

\For{$i = 1, \ldots, 4$}{
  $K_\mu^+ \gets$ solve barrier DARE at
    $\alpha_0 + \varepsilon\, e_i$\\
  $K_\mu^- \gets$ solve barrier DARE at
    $\alpha_0 - \varepsilon\, e_i$\\
  $\partial K_\mu / \partial \alpha_i
    \gets (K_\mu^+ - K_\mu^-) / (2\varepsilon)$
  \tcp*{central diff}
}
Store $(K_{\mu,0}, \partial_1 K_\mu, \ldots, \partial_4 K_\mu, \alpha_0)$:
240 floats.

\medskip
\textbf{Online (per control step):}\\
$\alpha \gets$ current flat parameter\\
$\delta\alpha \gets \alpha - \alpha_0$\\
$(x^*, u^*) \gets \Phi(\alpha)$\\
$K_\mu \gets K_{\mu,0}
  + \sum_i \delta\alpha_i \, \partial_i K_\mu$
\tcp*{192 MACs}
$u \gets u^* - K_\mu \, (x - x^*)$
\tcp*{48 MACs}
$u \gets \clip(u, u_{\min}, u_{\max})$
\tcp*{4 comparisons}
\end{algorithm}

% ============================================================
\section{Foolproof Implementation}\label{sec:foolproof}
% ============================================================

Everything above reduces to the following.
No optimisation theory is needed to implement it.

\subsection{Constants (your quadrotor)}

Fix these once for your hardware:
\begin{center}
\renewcommand{\arraystretch}{1.2}
\begin{tabular}{@{}l l l@{}}
\toprule
\textbf{Symbol} & \textbf{Meaning} & \textbf{How to get it} \\
\midrule
$m$ & mass (kg) & weigh it \\
$g$ & gravity ($9.81$) & constant \\
$J_x, J_y, J_z$ & inertia (kg\,m$^2$)
  & CAD or swing test \\
$k_t$ & per-motor thrust coeff & motor datasheet \\
$k_m$ & per-motor torque coeff & motor datasheet \\
$\ell$ & arm length (m) & measure it \\
$Q \in \R^{12 \times 12}$ & state cost
  & diagonal, you choose \\
$R \in \R^{4 \times 4}$ & input cost
  & diagonal, you choose \\
$\Delta t$ & control period (s)
  & your timer \\
$u_{\min}, u_{\max}$ & actuator limits
  & motor datasheet \\
\bottomrule
\end{tabular}
\end{center}

\subsection{Offline (run once, on your laptop)}

\begin{algorithm}[H]
\SetAlgoLined
\caption{Offline Precomputation (Analytic or Barrier)}\label{alg:offline}
\KwIn{$m,\, g,\, J,\, k_t,\, k_m,\, \ell,\,
  Q,\, R,\, \Delta t$,\;
  (optional) $\mu, u_{\min}, u_{\max}, x_{\min}, x_{\max}$.}
\KwOut{$K_0 \in \R^{4 \times 12}$,\;
  $K_1, K_2, K_3, K_4 \in \R^{4 \times 12}$,\;
  $\alpha_0 \in \R^4$.}

\medskip

\textbf{1.\ Hover equilibrium.}\\
$\alpha_0 \gets (0,\; 0,\; 0,\; 0)$
\tcp*{at hover, $a = (0,0,g)$ but code uses $(0,0,0,0)$}
$u_0^* \gets (T^*/(4k_t)) \cdot \mathbf{1} \in \R^4$
where $T^* = mg$\\

\medskip
\textbf{2.\ Hover linearisation.}\\
$(A_0, B_0) \gets \texttt{linearise}(\alpha_0)$
\tcp*{Algorithm~\ref{alg:linearise}}

\medskip
\textbf{3.\ Solve DARE at hover.}\\
\If{Barrier variant ($\mu > 0$)}{
  $H_\mu \gets$ Eq.~\eqref{eq:barrier_hess} from $(u_0^*, u_{\min}, u_{\max})$\\
  $H_{x,\mu} \gets$ Eq.~\eqref{eq:barrier_hess_x} from $(x_0^*, x_{\min}, x_{\max})$\\
  $R_\mu \gets R + H_\mu$;\;
  $Q_\mu \gets \texttt{dlqr}(A_0, B_0, Q, R_\mu) + H_{x,\mu}$\\
  $P_0 \gets \texttt{dare}(A_0, B_0, Q_\mu + \rho I, R_\mu + \rho I)$\\
}{
  $P_{\text{lqr}} \gets \texttt{dlqr}(A_0, B_0, Q, R)$\\
  $P_0 \gets \texttt{dare}(A_0, B_0, P_{\text{lqr}} + \rho I, R + \rho I)$\\
}
$K_0 \gets (R_{\text{aug}} + B_0^\trans P_0\, B_0)^{-1}\,
  B_0^\trans P_0\, A_0$\\

\medskip
\textbf{4.\ Central differences (4 directions).}\\
$\varepsilon \gets 10^{-3}$\\
\For{$i = 1, \ldots, 4$}{
  $\alpha^+ \gets \alpha_0 + \varepsilon\, e_i$,\quad
  $\alpha^- \gets \alpha_0 - \varepsilon\, e_i$\\
  $(A^{\pm}, B^{\pm}) \gets \texttt{linearise}(\alpha^{\pm})$\\
  $K^{\pm} \gets$ solve DARE (same procedure as Step 3)
    at $\alpha^{\pm}$\\
  $K_i \gets (K^+ - K^-)/(2\varepsilon)$
  \tcp*{central difference}
}

\medskip
\textbf{5.\ Store.}\\
Save $(K_0, K_1, K_2, K_3, K_4, \alpha_0)$:
$5 \times 4 \times 12 = 240$ numbers.\\
\end{algorithm}

\subsection{$\texttt{linearise}(\alpha)$}

\begin{algorithm}[H]
\SetAlgoLined
\caption{$\texttt{linearise}(\alpha)$:
  linearisation at flat parameter $\alpha$}
\label{alg:linearise}
\KwIn{$\alpha = (a_x, a_y, a_z, \psi)$,\;
  quadrotor constants $m, g, k_t, k_m, \ell, J, \Delta t$.}
\KwOut{$(A_d, B_d) \in \R^{12\times 12}
  \times \R^{12\times 4}$.}

\medskip
\tcp{Equilibrium attitude from flat output}
$T^* \gets m\,\sqrt{a_x^2 + a_y^2 + (a_z+g)^2}$\\
$\phi^* \gets \arcsin\bigl(
  m\,(a_x \sin\psi - a_y \cos\psi)\,/\,T^*
\bigr)$\\
$\theta^* \gets \arctan\bigl(
  (a_x \cos\psi + a_y \sin\psi)\,/\,(a_z+g)
\bigr)$\\

\medskip
\tcp{Gravity coupling matrix $G \in \R^{3\times 3}$}
$(s_\phi, c_\phi, s_\theta, c_\theta, s_\psi, c_\psi)
\gets$ trig of $(\phi^*, \theta^*, \psi)$\\
$G[0,:] \gets T^*/m \cdot
  (-c_\psi s_\theta s_\phi + s_\psi c_\phi,\;
   c_\psi c_\theta c_\phi,\;
   -s_\psi s_\theta c_\phi + c_\psi s_\phi)$\\
$G[1,:] \gets T^*/m \cdot
  (-s_\psi s_\theta s_\phi - c_\psi c_\phi,\;
   s_\psi c_\theta c_\phi,\;
   c_\psi s_\theta c_\phi + s_\psi s_\phi)$\\
$G[2,:] \gets T^*/m \cdot
  (-c_\theta s_\phi,\;
   -s_\theta c_\phi,\;
   0)$\\

\medskip
\tcp{Assemble $A_c$, $B_c$ (continuous time)}
$A_c[0{:}3, 6{:}9] \gets I_3$;\;
$A_c[3{:}6, 9{:}12] \gets W(\phi^*, \theta^*)$;\;
$A_c[6{:}9, 3{:}6] \gets G$\\
$B_c[6{:}9, :] \gets k_t/m \cdot e_3^{\mathcal{B}} \cdot \mathbf{1}^\trans$;\;
$B_c[9{:}12, :] \gets J^{-1} M_{\text{mix}}$\\

\medskip
\tcp{Discretise (forward Euler)}
$A_d \gets I_{12} + A_c\,\Delta t$;\;
$B_d \gets B_c\,\Delta t$\\
\Return $(A_d, B_d)$
\end{algorithm}

\subsection{Online (every control cycle)}

\begin{algorithm}[H]
\SetAlgoLined
\caption{Online Control Loop (the whole thing)}\label{alg:online}
\KwIn{Stored constants
  $K_0, K_1, K_2, K_3, K_4 \in \R^{4\times 12}$,\;
  $\alpha_0 = (0,0,0,0)$,\;
  $u_{\min}, u_{\max} \in \R^4$,\;
  mass $m$, gravity $g$, thrust coeff $k_t$.}

\medskip
\tcp{Run this at every time step}
\While{true}{

  \textbf{A.\ Read sensors.}\\
  $x \gets$ IMU + estimator output
  $= (p_x, p_y, p_z,\;
  \phi, \theta, \psi,\;
  v_x, v_y, v_z,\;
  \omega_x, \omega_y, \omega_z)^\trans
  \in \R^{12}$\\

  \medskip
  \textbf{B.\ Get flat parameter.}\\
  $(a_x, a_y, a_z, \psi_{\text{ref}}) \gets$
  trajectory planner output
  \tcp*{commanded accel + yaw}
  $\alpha \gets (a_x,\; a_y,\; a_z,\;
    \psi_{\text{ref}})$\\
  $\delta\alpha \gets \alpha - \alpha_0$
  \tcp*{4 subtractions}

  \medskip
  \textbf{C.\ Flat equilibrium.}\\
  $T^* \gets m\,\sqrt{a_x^2 + a_y^2 + (a_z+g)^2}$\\
  $\phi^* \gets \arcsin\bigl(
    m\,(a_x\sin\psi_{\text{ref}}
    - a_y\cos\psi_{\text{ref}})/T^*\bigr)$\\
  $\theta^* \gets \arctan\bigl(
    (a_x\cos\psi_{\text{ref}}
    + a_y\sin\psi_{\text{ref}})/(a_z+g)\bigr)$\\
  $x^* \gets (p^{\text{ref}},\;
    \phi^*, \theta^*, \psi_{\text{ref}},\;
    v^{\text{ref}},\;
    0, 0, 0)^\trans$\\
  $u^* \gets (T^*/(4k_t)) \cdot \mathbf{1}$\\

  \medskip
  \textbf{D.\ Form gain (4 scalar--matrix multiplies
    + 4 adds).}\\
  $K \gets K_0$
  \tcp*{copy $4 \times 12 = 48$ numbers}
  \For{$i = 1,\ldots,4$}{
    $K \gets K + \delta\alpha_i \cdot K_i$
    \tcp*{48 MACs each, 192 total}
  }

  \medskip
  \textbf{E.\ Compute control
    (one $4 \times 12$ matvec).}\\
  $\tilde x \gets x - x^*$
  \tcp*{12 subtractions}
  $\tilde u \gets K \cdot \tilde x$
  \tcp*{48 MACs}
  $u \gets u^* - \tilde u$
  \tcp*{4 subtractions}

  \medskip
  \textbf{F.\ Clamp.}\\
  \For{$j = 1,\ldots,4$}{
    $u_j \gets \max(u_{\min,j},\;
      \min(u_{\max,j},\; u_j))$
  }

  \medskip
  \textbf{G.\ Send.}\\
  Write $u$ to motor controller.\\
  Wait for next timer tick.\\
}
\end{algorithm}

\begin{remark}[Operation count]
Steps A, B: sensor read + 4 subtractions.
Step C: 1 sqrt, 1 asin, 1 atan, 4 trig, 6 multiplies
($\approx 20$ ops on CORDIC, $\approx 50$ on FPU).
Step D: 192 MACs.
Step E: 48 MACs + 16 adds.
Step F: 4 comparisons.
Total: $\approx 260$ arithmetic operations
per control cycle.
At 1\,GFLOP/s (low-end MCU): 0.26\,$\mu$s.
At 200\,MHz FPGA with 48 parallel MACs:
$\sim$75\,ns.
\end{remark}

% ============================================================
\section{FPGA Architecture}\label{sec:fpga}
% ============================================================

The zero-iteration law \eqref{eq:brutal}
maps to hardware with no control flow:
\textbf{pure datapath}.

\subsection{Pipeline stages}

\begin{center}
\renewcommand{\arraystretch}{1.3}
\begin{tabular}{@{}c l l l@{}}
\toprule
\textbf{Stage} & \textbf{Operation}
  & \textbf{Hardware} & \textbf{Latency} \\
\midrule
1 & $\alpha \gets (\ddot\sigma, \psi)$
  & register read & 1 clk \\
2 & $\delta\alpha \gets \alpha - \alpha_0$
  & $d$ adders & 1 clk \\
3 & $x^*(\alpha), u^*(\alpha)$
  & CORDIC (trig) & $\sim$10 clk \\
4 & $\tilde x \gets x_0 - x^*$
  & $n$ adders & 1 clk \\
5 & $K(\alpha) = K_0 +
    \sum_i \delta\alpha_i K_i$
  & $d\!\cdot\! mn$ MACs & $\sim$\!$mn/P$ clk \\
6 & $\tilde u = K(\alpha)\,\tilde x$
  & $mn$ MACs & $\sim$\!$mn/P$ clk \\
7 & $u = u^* - \tilde u$;\; $\clip$
  & $m$ add + compare & 1 clk \\
\bottomrule
\end{tabular}
\end{center}

\noindent
$P$ = number of parallel MAC units.
With $P = 48$ (one MAC per element of $K$):
stages 5--6 each take 1 clock cycle.
Total latency: $\sim$15 clock cycles.
At 200\,MHz: $\sim$75\,ns.

\subsection{Resource estimate (quadrotor)}

\begin{center}
\renewcommand{\arraystretch}{1.3}
\begin{tabular}{@{}l r l@{}}
\toprule
\textbf{Resource} & \textbf{Count} & \textbf{Note} \\
\midrule
Constant matrices & $1 + d = 5$
  & $K_0, K_1, \ldots, K_4 \in \R^{4\times 12}$ \\
Stored coefficients & $5 \times 48 = 240$
  & float16 or fixed-point \\
Storage & 480 bytes & fits in BRAM \\
DSP slices & $\le 48$
  & one per $K$ entry \\
CORDIC cores & 3
  & $\sin, \cos, \arctan$ \\
\bottomrule
\end{tabular}
\end{center}

\noindent
Fits on a low-end FPGA
(e.g.\ Xilinx Artix-7, Lattice iCE40).
No soft-core CPU, no OS, no memory allocator.
The controller is a circuit.

% ============================================================
\section{Complete Hierarchy}\label{sec:hierarchy}
% ============================================================

The five variants form a complexity cascade:
\begin{equation*}
\text{Online Relin}
\xrightarrow{\text{cache}}
\text{Flat Template}
\xrightarrow{K\to 0}
\text{Barrier Templ.}
\xrightarrow{\mu\to 0}
\text{Analytic Templ.}
\xrightarrow{\text{1st order}}
\text{Gain Sched.}
\end{equation*}

\begin{center}
\renewcommand{\arraystretch}{1.3}
\begin{tabular}{@{}l c c c c c@{}}
\toprule
\textbf{Variant}
  & \textbf{Linearisation}
  & \textbf{Iters}
  & \textbf{Online cost}
  & \textbf{Constraints}
  & \textbf{Valid region} \\
\midrule
Online Relin
  & fresh Jacobian
  & $K$
  & $O(n^3 + Kn^2)$
  & Hard
  & global \\
Flat Template
  & $\alpha$-grid cache
  & $K$
  & $K \cdot O(n^2)$
  & Hard
  & flat envelope \\
Barrier Templ.
  & hover + $\partial K_\mu$
  & 0
  & $O(d\,mn)$
  & Soft
  & flat envelope \\
Analytic Templ.
  & hover + $\partial K$
  & 0
  & $O(d\,mn)$
  & Clip
  & flat envelope \\
Gain Schedule
  & hover (fixed)
  & 0
  & $O(d\,mn)$
  & Clip
  & near hover \\
\bottomrule
\end{tabular}
\end{center}

\noindent
All variants with $\alpha$-parameterization share the
uniform tunnel guarantee (Theorem~\ref{thm:uniform}).
The key insight: flatness compresses the scheduling space
from $\R^{12}$ to $\R^4$ via
$\alpha = (a_x, a_y, a_z, \psi)$,
so only 4 sensitivity matrices (240 floats) are needed.

% ============================================================
\section{Experiment Results}\label{sec:results}
% ============================================================

All simulations use a Crazyflie 2.1 model
(mass 35\,g, arm 46\,mm) at 50\,Hz.
Python reference implementation; wall times are
for relative comparison only.

\subsection{Hover stabilization}

Initial offset $(0.2, 0.2, -0.2)$\,m.
$N=10$, $\rho=85$, 100 steps.

\begin{center}
\renewcommand{\arraystretch}{1.2}
\begin{tabular}{@{}lcccc@{}}
\toprule
 & TinyMPC & Online Relin & Flat Templ. & Analytic \\
\midrule
Avg.\ pos.\ err (m) & \textbf{0.064} & 0.095 & 0.138 & 0.145 \\
Avg.\ ADMM iters     & 80.0 & 16.7 & 7.8 & 0 \\
Wall time (s)        & 2.85 & 12.13 & \textbf{0.44} & \textbf{0.13} \\
\bottomrule
\end{tabular}
\end{center}

At hover, TinyMPC's fixed linearization is exact
($\alpha_0$ is hover) and achieves the best error.
But the Analytic Template is $22\times$ faster.

\subsection{Figure-eight tracking}

$xz$-plane, amplitude 0.5\,m, period 3.7\,s.
$N=15$, $\rho=5$, 200 steps.

\begin{center}
\renewcommand{\arraystretch}{1.2}
\begin{tabular}{@{}lcccc@{}}
\toprule
 & TinyMPC & Online Relin & Flat Templ. & Analytic \\
\midrule
Avg.\ pos.\ err (m) & \textbf{0.033} & 0.035 & 0.035 & 0.036 \\
Avg.\ ADMM iters     & 8.0 & 5.1 & 3.6 & 0 \\
Wall time (s)        & 1.16 & 22.02 & 0.75 & \textbf{0.26} \\
\bottomrule
\end{tabular}
\end{center}

All four variants achieve comparable accuracy.
The Flat Template is $1.5\times$ faster than TinyMPC;
the Analytic Template is $4.5\times$ faster.
Online Relin is $29\times$ slower due to
per-step DARE solves.

\subsection{Circle-in-square (constrained tracking)}

Circular trajectory ($R = 0.3$\,m) in the $xy$-plane,
inscribed square $|x|, |y| \le L = R/\sqrt{2}$.
$N=20$, $\rho=50$, 500 steps.

\begin{center}
\renewcommand{\arraystretch}{1.2}
\begin{tabular}{@{}lcccc@{}}
\toprule
 & Online Relin & Flat Templ. & Barrier & Analytic \\
 & (ADMM) & (ADMM) & (zero-iter) & (zero-iter) \\
\midrule
Avg.\ pos.\ err (m) & 0.062 & 0.062 & \textbf{0.017} & \textbf{0.017} \\
Max viol.\ (mm)      & \textbf{0.0} & \textbf{0.0} & 88.4 & 88.4 \\
Mean viol.\ (mm)     & \textbf{0.0} & \textbf{0.0} & 45.8 & 45.7 \\
Avg.\ ADMM iters     & 450.8 & 450.8 & 0 & 0 \\
Wall time (s)        & 168.1 & 157.4 & \textbf{0.65} & \textbf{0.62} \\
\bottomrule
\end{tabular}
\end{center}

\textbf{Key observations:}
\begin{itemize}[leftmargin=6mm]
  \item ADMM variants achieve 0\,mm violation by encoding
    absolute constraints as per-step error-state bounds.
  \item Zero-iteration variants show 88.4\,mm violation
    $= R - L$ (geometric limit).
  \item Barrier $=$ Analytic because the barrier Hessian
    cannot override a reference outside the feasible region.
  \item $240\times$ wall-time speedup from ADMM to zero-iteration.
\end{itemize}

% ============================================================
\section{Python API Reference}\label{sec:api}
% ============================================================

\subsection{Source modules}

\begin{center}
\renewcommand{\arraystretch}{1.2}
\begin{tabular}{@{}l l@{}}
\toprule
\textbf{Module} & \textbf{Description} \\
\midrule
\texttt{src/quadrotor.py} & 13-state quaternion dynamics (autograd) \\
\texttt{src/tinympc.py} & Core ADMM-based MPC solver \\
\texttt{src/flat\_linearization.py} & $\alpha$-parameterized analytic Jacobians \\
\texttt{src/flat\_template.py} & Flat Template: DARE cache grid + ADMM \\
\texttt{src/analytic\_template.py} & Analytic Template: $K_0 + \sum \delta\alpha_i K_i$ \\
\texttt{src/barrier\_template.py} & Barrier Template: log-barrier in DARE cache \\
\texttt{src/rho\_adapter.py} & ADMM penalty $\rho$ adaptation \\
\bottomrule
\end{tabular}
\end{center}

\subsection{Flat linearization API}

\begin{lstlisting}[caption={flat\_linearization.py}]
from src.flat_linearization import (
    get_quad_params,    # QuadrotorDynamics -> dict
    flat_equilibrium,   # alpha -> (x_eq, u_eq)
    linearise,          # alpha -> (A_d, B_d)
    quat_to_euler_state,  # 13D quat -> 12D Euler
    euler_to_quat,      # (phi, theta, psi) -> quat
)

# Setup
quad = QuadrotorDynamics(...)
params = get_quad_params(quad)

# Flat parameter: [a_x, a_y, a_z, psi]
alpha = np.array([0.0, 0.0, 0.0, 0.0])  # hover

# Equilibrium state (12D) and input (4D motor thrusts)
x_eq, u_eq = flat_equilibrium(alpha, params)

# Analytic Jacobians (no autograd)
A_d, B_d = linearise(alpha, params)  # (12x12), (12x4)
\end{lstlisting}

\subsection{Analytic Template API}

\begin{lstlisting}[caption={analytic\_template.py}]
from src.analytic_template import AnalyticTemplate

ctrl = AnalyticTemplate(params, Q, R, rho=5.0)
ctrl.precompute()  # solves 9 DAREs, stores 240 floats

# Online: ~260 ops per call
alpha = np.array([a_x, a_y, a_z, psi])
u = ctrl.control(x, alpha, u_min=u_lo, u_max=u_hi)
\end{lstlisting}

\subsection{Barrier Template API}

\begin{lstlisting}[caption={barrier\_template.py}]
from src.barrier_template import BarrierTemplate

ctrl = BarrierTemplate(
    params, Q, R, rho=50.0, mu=1e-3,
    u_min=u_lo, u_max=u_hi,    # input constraints
    x_min=x_lo, x_max=x_hi,    # state constraints
)
ctrl.precompute()  # barrier-augmented DAREs

# Online: ~260 ops (identical datapath)
u = ctrl.control(x, alpha, u_min=u_lo, u_max=u_hi)
\end{lstlisting}

\subsection{Flat Template API}

\begin{lstlisting}[caption={flat\_template.py}]
from src.flat_template import FlatTemplateMPC

# Grid of alpha values for cache
grid = [np.array([0, 0, 0, 0]), ...]

ctrl = FlatTemplateMPC(params, Q, R, N=15, rho=5.0,
                       grid_alphas=grid)
ctrl.set_bounds(umin=u_lo, umax=u_hi,
                xmin=x_lo, xmax=x_hi)
ctrl.set_tols_iters(max_iter=10,
                    abs_pri_tol=1e-3, abs_dua_tol=1e-3)

# Online: ADMM solve with cached matrices
x_traj, u_traj, status, iters = ctrl.solve(
    x0_error, alpha
)
\end{lstlisting}

\subsection{TinyMPC API}

\begin{lstlisting}[caption={tinympc.py}]
from src.tinympc import TinyMPC

ctrl = TinyMPC(A_hover, B_hover, Q, R,
               Nsteps=15, rho=5.0)
ctrl.set_bounds(umin=u_lo, umax=u_hi,
                xmin=x_lo, xmax=x_hi)

# Online: ADMM with fixed (A, B)
x_traj, u_traj, status, iters = ctrl.solve_admm(
    x_init, u_init, x_ref=x_ref, u_ref=u_ref
)
\end{lstlisting}

% ============================================================
\section{Discussion}\label{sec:discussion}
% ============================================================

\subsection{Connection to TTT}

The flat template is the LQ envelope
of \cite{zhang2026ttt} parameterised by
the flat coordinate $\alpha$.
The envelope moves with the operating point;
flatness provides the global chart
in which this motion is smooth.
The tunnel theorem of \cite{zhang2026ttt}
gives the contraction guarantee at each $\alpha$;
Theorem~\ref{thm:uniform} is the uniform version
over all $\alpha$.
The analytic template
(Proposition~\ref{prop:analytic-gain})
is the first-order Taylor expansion of the
envelope in $\alpha$---the tunnel
linearised in its own parameter space.

\subsection{Why Barrier Cannot Enforce Absolute Constraints}

The barrier Hessian penalizes deviations from the
\emph{reference} in error coordinates.
When the reference itself violates the constraint
(e.g., a circular trajectory outside a square bound),
no error-state cost modification can prevent violation.
The costate $\nabla V$ at the reference does not capture
the directional constraint information needed near boundaries.
Only ADMM's iterative projection can map absolute bounds
to per-step error-state bounds---the ``missing ingredient''
that zero-iteration controllers lack.

\begin{remark}[The black hole analogy --- 2026-02-18]
The constraint boundary acts like an event horizon:
the barrier Hessian $H_{x,\mu}$ grows as
$1/d^2$ near the boundary (where $d$ is the
distance to the constraint wall),
curving the cost landscape exactly like
the Schwarzschild metric curves spacetime near $r_s$.
The zero-iteration controller ``sees'' this curvature
at the reference point, but if the reference
is already past the horizon ($R > L$), no local
cost modification can pull the trajectory back---just
as no local force can escape a black hole.
The ADMM projection is the ``rocket'' that can
cross the horizon boundary by operating in
absolute coordinates.
We reproduced this black hole in the circle-in-square
experiment: 88.4\,mm of inescapable violation.
\end{remark}

\subsection{Gain scheduling, done right}

Classical gain scheduling computes
$K(\alpha)$ by interpolating
between hand-tuned operating points.
The analytic template produces
the same object but:
\begin{enumerate}[leftmargin=6mm]
  \item The schedule variable $\alpha$
    is \emph{derived} (from flatness),
    not chosen by intuition.
  \item The gain perturbation $K_i$
    is \emph{computed} (from the DARE),
    not tuned.
  \item The valid region is
    \emph{guaranteed} (by the tunnel),
    not assumed.
\end{enumerate}

\subsection{Higher-order templates}

The first-order expansion \eqref{eq:K-analytic}
is accurate near hover.
For extreme manoeuvres,
extend to second order:
$K(\alpha) = K_0 + \sum_i \delta\alpha_i K_i
+ \sum_{i,j} \delta\alpha_i \delta\alpha_j K_{ij}
+ \cdots$
Each $K_{ij}$ is a constant matrix.
Online cost scales as $O(d^2 mn)$---still
$O(1)$ iterations, just more MACs.
On FPGA, second order adds $\binom{d+1}{2} = 10$
additional scalar--matrix multiplies for $d = 4$.

\subsection{Beyond quadrotors}

Any differentially flat system admits
an analytic template:
\begin{itemize}[leftmargin=6mm]
  \item \textbf{Fixed-wing aircraft}:
    $\sigma = (x, y, z, \psi)$, $d = 4$.
  \item \textbf{Planar bipedal robot}:
    flat output = CoM + swing foot.
  \item \textbf{Differentially flat manipulators}:
    $\sigma$ = joint positions (trivially flat),
    $d = m$.
    Degenerates to computed-torque controller.
\end{itemize}

\subsection{What this does not do}

The analytic template addresses model mismatch
from operating-point dependence.
It does \emph{not} address:
\begin{itemize}[leftmargin=6mm]
  \item Unmodelled dynamics (wind, flex, latency)
    --- use a disturbance observer or L1 adaptive.
  \item Non-flat systems --- no flat coordinate
    chart exists; use the grid-based template.
  \item Hard constraint satisfaction under
    sustained saturation --- fall back to ADMM.
\end{itemize}

% ============================================================
% References
% ============================================================
\bibliographystyle{plain}

\begin{thebibliography}{9}

\bibitem{tinympc}
K.~Nguyen, S.~Schoedel, A.~Alavilli, B.~Plancher,
and Z.~Manchester.
\newblock {TinyMPC}: Model-predictive control on
  resource-constrained microcontrollers.
\newblock \emph{arXiv:2310.16985}, 2023.

\bibitem{zhang2026ttt}
J.~Zhang.
\newblock The tunnel theory.
\newblock 2026.

\bibitem{lancaster1995algebraic}
P.~Lancaster and L.~Rodman.
\newblock \emph{Algebraic Riccati Equations}.
\newblock Oxford University Press, 1995.

\bibitem{boyd2004}
S.~Boyd and L.~Vandenberghe.
\newblock \emph{Convex Optimization}.
\newblock Cambridge University Press, 2004.

\bibitem{jordana2023}
A.~Jordana, S.~Kleff, A.~Meduri, J.~Carpentier,
N.~Mansard, and L.~Righetti.
\newblock Stagewise {N}ewton method for dynamic game
  control problems.
\newblock \emph{arXiv:2307.09298}, 2023.

\bibitem{rugh2000}
W.~Rugh and J.~Shamma.
\newblock Research on gain scheduling.
\newblock \emph{Automatica}, 36(10):1401--1425, 2000.

\bibitem{hewer1971}
G.~Hewer.
\newblock An iterative technique for the computation
  of the steady-state gains for the discrete optimal
  regulator.
\newblock \emph{IEEE Trans.\ Automat.\ Control},
  16(4):382--384, 1971.

\bibitem{konstantinov1993}
M.~Konstantinov, P.~Petkov, and N.~Christov.
\newblock Perturbation analysis of the discrete
  Riccati equation.
\newblock \emph{Kybernetika}, 29(1):18--29, 1993.

\end{thebibliography}

\end{document}
